% Generated by tools/build_v6.py; edit the assembly rules, not this file.
\section{Convex comparison, common pairs, and additive localization}
\label{sec:transfer-foundations}

This appendix proves the specialized transfer inputs used by the scalar
lower bound and the stopping-program comparison. The finite-experiment
convex-estimation principle has antecedents in
Juditsky--Nemirovski and Polyanskiy--Wu \cite{jn,pw};
the midpoint adaptive decomposition is due to Lee--Valiant \cite{leevaliant}.
The arguments needed for this model are included in full.

\subsection{Finite categorical convex comparison}

\begin{lemma}[Finite categorical convex comparison]
\label{lem:finite-comparison}
Let $X$ be a nonempty compact convex subset of a finite-dimensional real vector
space. Let $\ell:X\to\R$ be affine, and let $p(x)$ be a strictly positive
categorical probability vector with affine continuous coordinates. For $n\ge1$
independent observations from $p(x)$, define
\[
 \Omega_n=\max\{\ell(x)-\ell(y):x,y\in X,
                  \operatorname{Aff}(p(x),p(y))^n\ge1/2\}.
\]
For every $\eta>0$, there is a bounded additive estimator
$F=c+\sum_{i=1}^n\phi_i(Z_i)$ satisfying
\begin{equation}
 \sup_{x\in X}\{\E_x(F-\ell(x))^2\}^{1/2}
 \le\tfrac32\Omega_n+\eta.
 \label{eq:finite-comparison}
\end{equation}
\end{lemma}

\begin{proof}
We reproduce the short argument needed for the numerical constant. Set
$s=\log2$, let $\mathcal F$ be the finite-dimensional space of additive
functions, and define, for $\alpha>0$,
\[
 \Phi(x,y;\phi,\alpha)=\ell(x)-\ell(y)
 +\alpha\log\E_y e^{\phi/\alpha}
 +\alpha\log\E_x e^{-\phi/\alpha}+2\alpha s.
\]
The function is continuous and concave in $(x,y)$: independence writes each
log moment as a sum of logarithms of positive affine forms. It is convex in
$(\phi,\alpha)$ by the perspective transformation of a log moment generating
function. Sion's minimax theorem \cite[corollary following Theorem 3.4, p.~174]{sion} therefore applies to $X\times X$ and
$\mathcal F\times(0,\infty)$. Only $X\times X$ needs to be compact;
the open, unbounded $\alpha$ domain is convex and causes no problem in
Sion's one-compact-set form. The infimum in $(\phi,\alpha)$ need not be
attained, and the later argument uses a strict slack $R>S$. Write
$2S=\inf_{\phi,\alpha>0}\max_{x,y}\Phi$.

For a fixed pair, the half log likelihood ratio belongs to $\mathcal F$ and
Cauchy--Schwarz gives
\[
 \inf_{\psi\in\mathcal F}
 \bigl[\log\E_y e^\psi+\log\E_x e^{-\psi}\bigr]
 =2\log\operatorname{Aff}(p(x)^{\otimes n},p(y)^{\otimes n}).
\]
Substituting $\phi=\alpha\psi$ and minimizing over $\alpha>0$ gives the value
$\ell(x)-\ell(y)$ when the affinity is at least $1/2$, and $-\infty$ otherwise.
Consequently $2S=\Omega_n$. The constrained maximum exists because its feasible
set is nonempty and compact.

For $R>S$, choose $(\phi,\alpha)$ with $\max\Phi\le2R$. On the diagonal,
Cauchy--Schwarz implies $\Phi(x,x;\phi,\alpha)\ge2\alpha s$, hence
$\alpha\le R/s$. Set
\[
 U=\max_x\{\ell(x)+\alpha\log\E_xe^{-\phi/\alpha}\},\qquad
 V=\max_x\{-\ell(x)+\alpha\log\E_xe^{\phi/\alpha}\}.
\]
Then $U+V+2\alpha s\le2R$. The estimator
$F=\phi+(U-V)/2$ satisfies
\[
 \E_xe^{\pm(F-\ell(x))/\alpha}\le e^{R/\alpha-s}.
\]
Therefore $\Pp_x\{|F-\ell(x)|>R+t\}\le e^{-t/\alpha}$ for $t\ge0$.
Tail integration yields
\[
 \E_x(F-\ell(x))^2
 \le R^2+2R\alpha+2\alpha^2
 \le R^2\left(1+\frac2{\log2}+\frac2{(\log2)^2}\right)<9R^2.
\]
Let $R\downarrow S$, with the slack chosen for $\eta$. Every component is
bounded because its alphabet is finite. This proves
\eqref{eq:finite-comparison}.
\end{proof}


\subsection{Midpoint decomposition for interleaved observations}

We use $\mathfrak h(P,Q)=1-\Aff(P,Q)$, so $0\le\mathfrak h\le1$ and
$\operatorname{TV}(P,Q)\le\sqrt{2\mathfrak h(P,Q)}$.
The following finite-prefix argument is the midpoint decomposition of
Lee--Valiant \cite[Lemmas 17--18]{leevaliant}, written in the form needed here.

\begin{lemma}[Midpoint work decomposition]\label{lem:tf-midpoint}
Fix two endpoint priors $q_-,q_+$ and their midpoint
$\bar q=(q_-+q_+)/2$. Suppose every capped randomized one-unit policy
$A$ obeys
\[
 \mathfrak h(P_-^A,P_+^A)\le\kappa\,\E_{\bar q}T_A.
\]
Then every population algorithm with a deterministic total cap satisfies
\[
 \mathfrak h(P_-^{\rm pop},P_+^{\rm pop})
 \le\kappa\,\E_{\bar q}T_{\rm pop}.
\]
It suffices that the local hypothesis hold up to that same total cap.
If instead $q_\pm\le1/2$ and the local bound uses null work
$\lambda\E_0T_A$, the conclusion holds with $\kappa=2\lambda$.
The statement also holds for finitely many observable component types,
each with its own endpoint pair and midpoint, when the same coefficient
bounds every local policy for every type. Types may be selected adaptively, but each fresh unit's type is fixed
before its first observation and retained on every revisit.
\end{lemma}

\begin{proof}
Fix total cap $K$ and condition on the algorithm's independent seed.
Defer latent draws until a unit is first observed, which changes no
observable law. Index units by that first observation. A consistent global leaf $z$
contains local binary words $h_i(z)$, with empty words for unused slots.
The endpoint leaf densities factor as
\[
 p_\pm^{\rm pop}(z)=\prod_{i=1}^K p_\pm(h_i(z)).
\]
The common policy-consistency indicator is implicit. Put
$m_i=(p_-(h_i)+p_+(h_i))/2=p_{\bar q}(h_i)$ and
$c_i=\sqrt{p_-(h_i)p_+(h_i)}$. Then
\[
 \prod_i m_i-\prod_i c_i
 \le\sum_i(m_i-c_i)\prod_{j\ne i}m_j.
\]
The integral of $\prod_i m_i$ over consistent leaves is one, since it
is the population law in which every fresh unit has midpoint prior.
Integrating the inequality therefore bounds population squared Hellinger
by a sum of hybrid squared Hellinger distances. In hybrid $i$ only
unit $i$ has an endpoint prior; all other units have midpoint priors.

Simulate all other units from independent midpoint tapes included in
an auxiliary seed, and run the population rule with one genuine unit
in slot $i$. This is a capped local policy. Its full transcript,
including that seed, determines the hybrid global transcript.
Data processing bounds hybrid distance by local distance. The expected
midpoint local works sum to $\E_{\bar q}T_{\rm pop}$ because each
purchased observation belongs to exactly one slot. This proves the
first assertion, including after integrating the original seed.
The charged expectation is under the midpoint population, admissible
by convexity. It is not replaced by the average endpoint work, which
can be much smaller under an adaptive stopping rule.

If $q_\pm\le1/2$, the midpoint latent law dominates half the null
latent law. The same measure inequality holds after passing through
any observation policy and auxiliary seed. Thus
$\E_0T_A\le2\E_{\bar q}T_A$.

For observable types, the type selected for a slot is a function of the
history before its first observation. Use the resulting type-specific
kernel in each leaf product. To implement hybrid slot $i$, first generate
all observations before that slot under the componentwise midpoint law.
Conditional on this prehistory and its selected type, the genuine unit
has that type's prescribed endpoint law and is fresh, independently of
the remaining auxiliary randomness. This is conditional independence:
the latent law may depend on the type contained in the seed.

Retain the already generated tapes of other units on every revisit.
For later slots, an independent array of potential tapes indexed by slot
and type can be placed in the auxiliary seed; choose the appropriate
unused tape only after that slot's type has been selected. The seed's
distribution is the same under the two endpoint laws for the genuine
unit. Conditional on the prehistory, the local hypothesis for its selected
type applies. If slot $i$ is never reached, the genuine local work is zero.
Summing and integrating counts every purchased observation once under
the actual componentwise midpoint population. No latent magnitude or
sign of the genuine unit is revealed to the local policy.
\end{proof}

For an unbounded stopping rule, apply the conclusion to every bounded
global prefix. Total variation of the full transcript is the increasing
limit of the prefix total variations: bounded cylinder functions generate
the transcript sigma-field, or equivalently one applies conditional
expectation to the likelihood densities under the dominating sum law.
Finite expected endpoint work ensures almost-sure finite termination.
A common pair valid for every cap therefore gives the same terminal
testing bound. Equivalently, couple the stopped experiment with its
cap-$L$ truncation under each endpoint. Their total coupling error is at
most $\Pp_-(T>L)+\Pp_+(T>L)$, which tends to zero. This passage
requires no common deterministic bound on stopping.

For unbounded stopping in the budgeted applications, the same endpoint pair
works at every cap, and the original uniform budget gives expected work
at most $B$ at both endpoints and at their admissible midpoint. Let $N$
be original purchased work, and cap after $L$ observations, recording an
overflow flag if needed. Any terminal event restricted to $\{N\le L\}$
is measurable from this capped transcript. Therefore
\[
 \operatorname{TV}(P_-,P_+)
 \le\operatorname{TV}(P_-^{(L)},P_+^{(L)})
       +\Pp_-(N>L)+\Pp_+(N>L)
 \le\sqrt{2\kappa B}+2B/L.
\]
Letting $L\to\infty$ gives the terminal testing bound. For an oracle
simulation, cap at the original observation count; its genuine work
is no larger, and the same argument applies. This avoids exchanging a
cap limit with a separately chosen pair.

\subsection{The half-prior class and null work}

Put
\[
 X_\rho=\left\{q:0\le q\le\tfrac12,\quad
                  \theta(q):=\int q\,d\mu\le\rho\right\},
 \qquad 0<\rho\le\tfrac14.
\]
For a capped observable one-coin procedure with bounded real output \(Y\),
write
\[
 V=\sup_{q\in X_\rho}\Var_qY,\qquad
 M=\sup_{q\in X_\rho}\E_qY^2,\qquad W_0=\E_0T,
\]
and define
\begin{equation}
 \mathfrak J_0(b)=\inf\left\{VW_0:
       \sup_{q\in X_\rho}|\E_qY-\theta(q)|\le b,
       \ T\text{ capped},\ Y\text{ bounded}\right\}.
 \label{eq:null-bias-complexity}
\end{equation}
An empty infimum is infinity.  The function is nonincreasing in \(b\),
and \(\mathfrak J_0(\rho/2)=0\).

\subsection{Transfer at an arbitrary bias tolerance}

\begin{theorem}[Bias-indexed common pairs and population lower bounds]
\label{thm:bias-transfer}
Suppose that every capped bounded one-coin output with uniform bias at
most \(b>0\) over \(X_\rho\) satisfies \(VW_0\ge J\), where
\(0<J<\infty\).  Then, for every \(0<\delta\le b/3\), there is one pair
\(q_-,q_+\in X_\rho\) with target separation at least \(\delta\) such
that simultaneously for all capped randomized local policies \(A\),
\begin{equation}
 \mathfrak h(P_{q_-}^A,P_{q_+}^A)
 \le42\frac{\delta^2}{J}\E_0T_A.
 \label{eq:bias-common-pair}
\end{equation}
Consequently,
\begin{equation}
 R_{\mu,\rho}(B)\ge\frac1{128}
                 \min\{b,\sqrt{J/B}\}.
 \label{eq:bias-population-transfer}
\end{equation}
The conclusion includes fully adaptive creation, revisiting, interleaving,
nonlinear outputs and unbounded stopping.
\end{theorem}

\begin{proof}
The premise forces \(b<\rho/2\), since the zero-work constant \(\rho/2\)
would otherwise contradict it.  First fix a finite strictly positive
categorical statistic of a capped policy of null work \(w>0\).
For \(n\) independent runs, let
\(F=c+\sum_{i=1}^nY_i\) be a bounded additive estimator of worst RMS
\(r\).  Put \(b_i=\E_0Y_i\), \(b_F=\E_0F\), and choose an independent
uniform index \(I\).  Running just its component and returning
\(Z=n(Y_I-b_I)+b_F\) preserves all population means at null work \(w\).

The laws under \(q\) and zero share a subprobability of mass
\(1-\theta(q)\), obtained by weighting the negative-sign kernel by
\(1-q(d)\).  The elementary common-mixture variance inequality gives
\[
 (\E_qY_i-b_i)^2\le
 \frac{2\theta(q)}{1-\theta(q)}
                 (\Var_qY_i+\Var_0Y_i).
\]
Indeed, either mean's displacement from the normalized common part is
at most the square root of \(\theta/(1-\theta)\) times its standard
deviation; square the sum of the two displacements.  Independence in
\(F\) now yields
\[
 \sup_q\Var_qZ\le
 nr^2\left(1+\frac{4\rho}{1-\rho}\right)\le\tfrac73nr^2.
\]
If \(r<b\), apply the assumed obstruction to \(Z\); if \(r\ge b\),
retain that inequality.  In both cases,
\[
 r\ge\min\{b,\sqrt{3J/(7nw)}\}.
\]

Lemma~\ref{lem:finite-comparison} gives an additive estimator of
RMS at most \((3/2)\Omega_n+\eta\), where \(\Omega_n\) is the maximum
target separation whose product affinity is at least \(1/2\).
Letting \(\eta\downarrow0\) therefore gives
\[
 \Omega_n\ge\tfrac23\min\{b,\sqrt{3J/(7nw)}\}.
\]
Set \(x=J/(21\delta^2w)\).  If \(x\ge2\), take
\(n=\lfloor x\rfloor\); then \(\Omega_n\ge2\delta\), and the
corresponding pair satisfies
\[
 \mathfrak h\le\frac{\log2}{n}
 \le42\log2\,\frac{\delta^2w}{J}
 \le42\frac{\delta^2w}{J}.
\]
If \(x<2\), the last right side exceeds one, so any separated pair
suffices.  If \(w=0\), the seed-only first-observation event has
probability zero; all local laws agree and \(q_-=0,q_+=\delta\)
suffices.

We now obtain one pair simultaneously for every capped local policy.
For an arbitrary capped policy, partition its full transcript into finitely
many measurable cells, then replace the cell independently by a uniform
cell with probability \(0<\varepsilon<1\).  Run the policy even on
replacement outcomes, so its null work is unchanged.  The resulting
strictly positive category probabilities and the target have compact
convex finite-dimensional range, so the preceding comparison applies.
For any finite collection of these experiments, independently mix them
and reveal the chosen index. Categories with zero mixture weight can be
dropped, leaving a strictly positive finite alphabet. For a fixed one-unit
policy, its categorical law is affine in $q$: condition on the one genuine
latent state first. This affineness is not being asserted for an entire
adaptive population transcript.

Set $k=42\delta^2/J$ and use the oriented, weak-star compact convex pair set
\[
 K_\delta=\{(q_-,q_+)\in X_\rho^2:
                 \theta(q_+)-\theta(q_-)\ge\delta\}.
\]
For experiment $i$, let $w_i$ be its null work and
$f_i(q_-,q_+)=\mathfrak h(P_{q_-}^{A_i},P_{q_+}^{A_i})-kw_i$.
Each $f_i$ is continuous and convex on $K_\delta$.
For any mixture weights $\lambda$, the revealed index gives exactly
$\mathfrak h_{\rm mix}=\sum_i\lambda_i\mathfrak h_i$ and $w_{\rm mix}=\sum_i\lambda_iw_i$.
The preceding one-experiment argument therefore implies
\[
 \min_{z\in K_\delta}\sum_i\lambda_i f_i(z)\le0
 \qquad\text{for every }\lambda\text{ in the finite simplex}.
\]
Applying convex minimax in its indicated direction gives
\[
 \min_{z\in K_\delta}\max_i f_i(z)
 =\max_\lambda\min_{z\in K_\delta}\sum_i\lambda_i f_i(z)\le0.
\]
The left minimum is attained. Thus the finite family has a single feasible
pair; separate pairs for the individual experiments would not suffice.
All its constraint sets are closed in $K_\delta$.
The finite intersection property supplies one pair for all experiments.
For fixed $b$ and $\delta$, this pair is chosen before the adaptive
population algorithm and its total-observation cap. The same pair
works for every capped local policy and every cap.
The intersection ranges over all caps, finite partitions and positive
noise levels at once, so this pair does not depend on the stopping cap.
The pair is fixed while $\varepsilon\downarrow0$; finite categorical
Hellinger distance is continuous even when some limiting cells have zero
mass. Full squared Hellinger distance is the supremum over finite
measurable partitions. For completeness, put $R=P+Q$ and $t=dP/dR$.
Partition the range of $t$ into intervals of width at most $1/m$.
Since $g(t)=\sqrt{t(1-t)}$ satisfies $|g(u)-g(v)|\le\sqrt{|u-v|}$,
the full and coarsened affinities differ by at most $2/\sqrt m$.
Thus the supremum statement holds without a countable-generation
assumption on the transcript space. This proves
\eqref{eq:bias-common-pair} for arbitrary transcript spaces.
The quantifier is one pair for every policy and cap at fixed
$(\mu,\rho,b,J,\delta)$. The pair may change when $\delta$, and hence
the chosen budget-dependent separation, changes.

The midpoint decomposition in Lemma~\ref{lem:tf-midpoint}
bounds population squared Hellinger by
\(84\delta^2B/J\): the local null cost is at most twice the midpoint
cost, and summing midpoint costs counts each observation once.  Set
\[
 \delta=\frac1{32}\min\{b,\sqrt{J/B}\}.
\]
The distance is at most \(84/1024<1/8\), hence TV is at most \(1/2\).
The same pair works for every deterministic observation cap.  Finite
expected endpoint work and increasing capped prefixes give the same
terminal TV bound for unbounded stopping.  Thresholding an estimator at
the target midpoint gives RMS at least \(\delta/4\), proving
\eqref{eq:bias-population-transfer}.  A budget over the full rarity class
implies the required budget over \(X_\rho\).
\end{proof}

\subsection{A direct variational null-work obstruction}

For Bernoulli observations, recall the variational modulus of
Section~\ref{sec:rare}
\begin{equation}
 \Psi_{\mu,\rho}(h)=
 \sup_{\substack{0\le w\le1\\\int w\,d\mu\le\rho}}
 \int\min\{w(d),h\sqrt{w(d)}/d\}\,\mu(dd).
 \label{eq:inherited-psi}
\end{equation}
The proof below establishes the exact null-work input required here
directly from fixed-sign comparisons. No optimizer formula is needed.

\begin{proposition}[Variational accuracy-dependent null obstruction]
\label{prop:psi-null}
For every probability measure \(\mu\) on \((0,1/2]\) and every
\(0\le b\le\rho/64\),
\begin{equation}
 \mathfrak J_0(b)\ge
 \sup_{h>0}\frac{(\Psi_{\mu,\rho}(h)-4b)_+^2}{192h^2}.
 \label{eq:psi-null-product}
\end{equation}
No nonatomicity, finite support, or separation of magnitudes is assumed.
\end{proposition}

\begin{proof}
Fix an eligible \((Y,T)\), and define its conditional response, moments,
and negative work by
\[
 f(d)=\E_{d,+}Y-\E_{d,-}Y,\qquad
 A_s(d)=\E_{d,s}Y^2,\qquad w_-(d)=\E_{d,-}T.
\]
We first show
\begin{equation}
 f(d)^2\le64d^2w_-(d)[A_-(d)+A_+(d)].
 \label{eq:conditional-response-cost}
\end{equation}
For \(d\le1/4\), the stopped KL chain rule gives
\(\KL(P_{d,-}^A,P_{d,+}^A)\le(64/3)d^2w_-(d)\).
For any square-integrable output under two laws,
\[
 |\E_PY-\E_QY|^2
 \le4(\E_PY^2+\E_QY^2)\mathfrak h(P,Q),\qquad
 \mathfrak h(P,Q)\le\tfrac12\KL(P,Q).
\]
These prove \eqref{eq:conditional-response-cost} with coefficient
\(128/3<64\).  For \(d>1/4\), the conditional laws agree whenever
the seed requests no observation.  Their Hellinger distance is therefore
at most the activation probability, hence at most \(w_-(d)\).
The same mean inequality gives
\(f(d)^2\le4w_-(d)(A_-+A_+)\le64d^2w_-(d)(A_-+A_+)\).
This also handles \(d=1/2\) without a singular endpoint KL comparison.

For every feasible \(w\) in \eqref{eq:inherited-psi}, the allocation
\(q=w/2\) belongs to \(X_\rho\), and
\[
 \int w(A_-+A_+)\,d\mu
 \le2\int[(1-w/2)A_-+(w/2)A_+]\,d\mu\le2M.
\]
Weighted Cauchy--Schwarz and \eqref{eq:conditional-response-cost} yield
\begin{equation}
 \left(\int\frac{\sqrt w\,|f|}{d}\,d\mu\right)^2
 \le W_0\int\frac{wf^2}{d^2w_-}\,d\mu
 \le128MW_0.
 \label{eq:weighted-response-bound}
\end{equation}
On \(\{w_-=0\}\), the response inequality forces \(f=0\), and the
quotient is zero.  Truncation justifies Cauchy--Schwarz before passing to
the limit if an integral is initially extended-valued.

Put \(v_h=\min\{w,h\sqrt w/d\}\).  Comparing the mean under the
admissible allocation \(q=v_h/2\) with the null mean, and using both
bias bounds, gives
\[
 \int v_h f\,d\mu\ge\int v_h\,d\mu-4b.
\]
Since \(v_h\le h\sqrt w/d\),
\eqref{eq:weighted-response-bound} implies
\[
 MW_0\ge\frac{(\int v_h\,d\mu-4b)_+^2}{128h^2}.
\]
Take the supremum over \(w\).  To convert the second moment to variance,
the null and constant-prior populations force the all-positive and
all-negative means to differ by at least \(31/32\).  The mixture with
constant sign prior \(\rho\) consequently has variance at least
\(\rho(1-\rho)(31/32)^2\ge2\rho/3\).  Also every admissible mean has
absolute value at most \(65\rho/64\).  Thus
\(M\le V+(65\rho/64)^2\le3V/2\), proving
\eqref{eq:psi-null-product}.
\end{proof}

This product bound accounts for accuracy-dependent tail difficulty in the
same framework as the annular scale-entropy obstruction.  At a fixed gap
\(d\), the directly evaluated identity
\(\Psi(h)=\min\{\rho,h\sqrt\rho/d\}\) recovers product scale
\(\rho/d^2\) as \(b\downarrow0\).  For the quadratic law it supplies
the inverse-square logarithm; the additional scale-entropy factor still
requires the annular argument.



The next proposition restates the optimizer of Section~\ref{sec:rare},
whose value formula is Lean-checked there; the short proof is retained
because the argument above uses its form.

\begin{proposition}[An optimizer for the variational modulus]
\label{prop:psi-optimizer}
For $0<\rho<1$, there is a unique $\alpha>0$ satisfying
\[
 \int\min\{1,\alpha^2/d^2\}\,\mu(dd)=\rho.
\]
The allocation $w_*(d)=\min\{1,\alpha^2/d^2\}$ maximizes
$\Psi_{\mu,\rho}(h)$ simultaneously for every $h>0$. In particular,
$\Psi(h)=\rho$ when $h\ge\alpha$. Fractional allocation at an atom is
permitted and may be essential; this is distinct from the hard
assignment statement for the entropy-selection infimum.
\end{proposition}
\begin{proof}
The mass on the left is continuous in $\alpha$, tends to zero at zero
and to one at infinity, and is strictly increasing while it is below
one. These assertions follow by dominated convergence and the positive
mass of unsaturated magnitudes. Thus the stated root exists uniquely.
If $h\ge\alpha$, the allocation attains the upper bound $\rho$.
If $h<\alpha$, put $\lambda=h/(2\alpha)<1/2$. The function
\[
 \phi_d(w)=\min\{w,(h/d)\sqrt w\},\qquad 0\le w\le1,
\]
is concave. For $d>\alpha$, the point $w_*=(\alpha/d)^2$ is in its
square-root branch and has derivative $\lambda$. For $h<d\le\alpha$,
the left derivative at $w_*=1$ is $h/(2d)\ge\lambda$; for $d\le h$,
$\phi_d(w)=w$. Hence $w_*$ maximizes $\phi_d(w)-\lambda w$
pointwise. Integrating and using its mass $\rho$ proves optimality over
all allocations of mass at most $\rho$, including atomic laws.
\end{proof}

\begin{remark}[The bias subtraction cannot be uniformly reduced]
Let $\rho=1/8$, $\varepsilon=2^{-10}$, and
$\mu=(1-\varepsilon)\delta_{1/2}+\varepsilon\delta_d$.
The one-flip output $Y=X-b$, with $b=\varepsilon(1-2d)/4$, has bias
at most $b$ on $X_\rho$, null work one, and variance at most $1/4$.
Indeed its bias is
$\varepsilon(1/2-d)(1-2q(d))-b\in[-b,b]$.
Yet $\Psi(d)\ge\varepsilon$. Replacing $4b$ by $cb$ with any fixed
$c<4$ in \eqref{eq:psi-null-product} would make its right-hand side
diverge as $d\downarrow0$. The coefficient four cannot be reduced in
this form of the bound, even if its denominator constant is changed.
\end{remark}

\subsection{A direct sum of conditional experiments}

Let $\sigma_1,\ldots,\sigma_m$ be nonzero finite submeasures with
$\sum_i\sigma_i\le\mu$. Write $\tau_i=\sigma_i((0,1/2])$ and
$\nu_i=\sigma_i/\tau_i$. Disjoint restrictions are a special case;
fractional restrictions of atoms are also allowed. Choose
\[
 0<r_i\le1/4,\qquad \sum_i\tau_i r_i\le\rho.
\]
Within component $i$, use the half-prior subclass
\[
 X_{r_i}(\nu_i)=\{0\le q\le1/2:\ \textstyle\int q\,d\nu_i\le r_i\}.
\]
The original population lower bound is valid on this subclass of the full
rarity class.

\begin{theorem}[Additive localization]\label{thm:additive-localization}
Suppose that, for each $i$, every capped bounded observable one-unit output
with uniform bias at most $b_i>0$ over $X_{r_i}(\nu_i)$ satisfies
\[
 \left(\sup_q\Var_qY\right)\E_{i,0}T\ge J_i>0.
\]
Then, for every $B>0$,
\begin{equation}\label{eq:additive-localization}
 \boxed{\quad
 R_{\mu,\rho}(B)\ge\frac1{128}
 \sum_{i=1}^m\tau_i\min\{b_i,\sqrt{J_i/B}\}.
 \quad}
\end{equation}
The same total budget $B$ occurs in every summand. The component label
is chosen before the first observation and retained on every revisit;
its conditional magnitude and sign remain hidden.
\end{theorem}

\begin{proof}
Theorem~\ref{thm:bias-transfer} applies within every component. Set
\[
 \delta_i=\frac1{32}\min\{b_i,\sqrt{J_i/B}\}.
\]
It supplies a pair $q_{i,-},q_{i,+}\in X_{r_i}(\nu_i)$, oriented to have
conditional target separation at least $\delta_i$, such that simultaneously
for every capped local policy $A$,
\begin{equation}\label{eq:additive-local-pair}
 \mathfrak h(P_{i,-}^A,P_{i,+}^A)
 \le42\frac{\delta_i^2}{J_i}\E_{i,0}T_A
 \le\frac{42}{1024B}\E_{i,0}T_A.
\end{equation}

Fix an arbitrary original algorithm with uniform expected budget $B$.
Consider an oracle experiment that can select a component at no cost and
then receive a genuine coin with magnitude law $\nu_i$. Its component label
is visible, while its magnitude and sign remain hidden. The algorithm can
adaptively select, revisit, and interleave these coins. Apply the argument
to the simulation of the fixed original algorithm described below. At the
componentwise midpoint it induces the original parameter
$\bar q=(q_-+q_+)/2\in\mathcal Q_\rho$, so its genuine midpoint work is
at most the original work and hence at most $B$. Under the conditional
midpoint prior, $q\le1/2$, so the midpoint law dominates half the null law.
Lemma~\ref{lem:tf-midpoint} therefore bounds terminal squared
Hellinger by $84/1024<1/8$: the local coefficient in
\eqref{eq:additive-local-pair} is common to all components, and the summed
conditional midpoint costs count each purchased observation once.
Increasing capped prefixes handles unbounded stopping. The total target
separation is at least $\sum_i\tau_i\delta_i$.

Every original algorithm can be simulated in this stronger experiment.
Privately choose the component label with probabilities $\tau_i$ whenever
it creates a fresh coin. On the remaining probability mass, simulate the
known all-negative residual law for free. Keep the chosen label when the
coin is revisited. For fractional submeasures, integrating this private
mixture label gives the original admissible parameters
\[
 q_\pm(d)=\sum_i\frac{d\sigma_i}{d\mu}(d)q_{i,\pm}(d),
 \qquad 0\le q_\pm\le1/2,\qquad
 \int q_\pm\,d\mu\le\sum_i\tau_i r_i\le\rho.
\]
The number of genuine oracle observations does not exceed the original
number. Thus the oracle lower bound applies to the original experiment.
The Hellinger bound gives TV at most $1/2$, and midpoint testing gives RMS
at least one quarter of the target separation. This is
\eqref{eq:additive-localization}.
\end{proof}

For a direct explicit input, suppose component $i$ is supported in
$(0,a_i]$, $a_i\le1/4$, and select a finite separated family of small-gap
annuli. Let $\Gamma_i^{\circ}$ and $w_{*,i}$ denote its conditional cost entropy
and minimum band mass. Proposition~\ref{prop:scaled-annular} gives
\[
 b_i=r_i/64,\qquad
 J_i=c_{\mathrm{ann}}
       \min\{r_i\Gamma_i^{\circ},w_{*,i}/a_i^2\},\qquad
 c_{\mathrm{ann}}=\frac{e^{-28}}{2{,}000{,}000}.
\]
Write the corresponding unnormalized quantities with respect to
$\sigma_i$ as $\Gamma_{\sigma_i}^{\circ}=\tau_i\Gamma_i^{\circ}$ and
$w_{*,\sigma_i}=\tau_iw_{*,i}$. The conversion uses the exact identity, with $\theta=\tau r$,
\[
 \tau^2\min\{r\Gamma^{\circ},w_*/a^2\}
 =\min\{\theta\Gamma_{\sigma}^{\circ},\tau w_{*,\sigma}/a^2\}.
\]
In particular the second branch retains the factor $\tau$. Consequently,
\begin{equation}\label{eq:additive-annular}
 R_{\mu,\rho}(B)\gtrsim
 \sum_i\min\left\{\tau_i r_i,
 \sqrt{\frac{
 \min\{\tau_i r_i\Gamma_{\sigma_i}^{\circ},
           \tau_iw_{*,\sigma_i}/a_i^2\}}B}\right\}.
\end{equation}
For clarity, the squared component mass is essential before cancellation:
\[
 \tau_i^2J_i=c_{\mathrm{ann}}
 \min\{\tau_i r_i\Gamma_{\sigma_i}^{\circ},
             \tau_iw_{*,\sigma_i}/a_i^2\}.
\]
This explains the extra $\tau_i$ in the saturation term of
\eqref{eq:additive-annular}; it must not be dropped when passing between
conditional and unnormalized quantities.
Supremizing this expression over finite decompositions and rarity
allocations gives an explicit lower functional. The bias-indexed form
\eqref{eq:additive-localization} also permits each component's accuracy
scale to be chosen separately. Appendix~\ref{sec:matching-lower} applies this expression to components
selected by the entropy--tail optimizer, and Appendix~\ref{fg:sec:matching-lower}
applies it to the components constructed by the quantile-core filter and
packing.


\subsection{A universal risk floor}

\begin{lemma}[Bernoulli-label floor]\label{lem:tf-floor}
For every magnitude law, $0<\rho\le1/4$, and $B>0$,
\[
 R_{\mu,\rho}(B)\ge\frac1{16}
       \min\{\rho,\sqrt{\rho/B}\}.
\]
\end{lemma}
\begin{proof}
Restrict to constant sign allocations. Give an oracle the exact sign
of a unit when its first observation is requested, for one charged
observation. It can sample $D\sim\mu$
independently and simulate all future flips for free, so every original
algorithm can be simulated with no greater work. The informative labels
are iid Bernoulli with the unknown constant prior.

Take $p_0=\rho/2$ and $p_1=p_0+\delta$, where
$\delta=\min\{\rho/2,\sqrt{\rho/B}/4\}$. Both are admissible and
$p_1(1-p_1)\ge3\rho/8$. The stopped KL chain rule, first at bounded
prefixes and then by monotone convergence, gives
\[
 \KL(P_0,P_1)\le B\,\operatorname{kl}(p_0,p_1)
 \le\frac{8B\delta^2}{3\rho}\le\frac16.
\]
Thus $\operatorname{TV}(P_0,P_1)\le1/2$. Thresholding any estimator at
the target midpoint implies worst RMS at least $\delta/4$.
Finally
$\delta/4=\min\{\rho/8,\sqrt{\rho/B}/16\}
\ge\min\{\rho,\sqrt{\rho/B}\}/16$.
\end{proof}
